% Part: first-order-logic
% Chapter: beyond
% Section: many-sorted-logic

\documentclass[../../../include/open-logic-section]{subfiles}

\begin{document}

\olfileid[pt]{fol}{byd}{msl}

\olsection{Lógica de Múltiplas Ordenações}

Na lógica de primeira ordem, variáveis e quantificadores variam sobre um único
!!{domain}. Mas é frequentemente útil ter vários !!{domain}s (disjuntos):
por exemplo, pode-se querer ter !!a{domain} de números,
!!a{domain} de objetos geométricos, !!a{domain} de funções de
números para números, !!a{domain} de grupos abelianos, e assim por diante.

A lógica de múltiplas ordenações fornece esse tipo de arcabouço. Começa-se com uma
lista de ``ordenações''---a ``ordenação'' de um objeto indica o
``!!{domain}'' que ele deveria habitar. Tem-se então !!{variable}s
e quantificadores para cada ordenação e (geralmente) uma !!{identity} para cada
ordenação. Funções e relações também são ``tipadas'' pelas ordenações dos
objetos que podem receber como argumentos. Caso contrário, mantêm-se as regras
usuais da lógica de primeira ordem, com versões das regras de quantificadores
repetidas para cada ordenação.

Por exemplo, para estudar relações internacionais, poderíamos escolher uma
linguagem com dois tipos de objetos, cidadãos franceses e cidadãos alemães.
Poderíamos ter uma relação unária, ``bebe vinho'', para objetos
da primeira ordenação; outra relação unária, ``come wurst'', para objetos
da segunda ordenação; e uma relação binária, ``forma um casal
multinacional'', que recebe dois argumentos, onde o primeiro argumento
é da primeira ordenação e o segundo argumento é da segunda ordenação. Se
usarmos variáveis $a$, $b$, $c$ para variar sobre cidadãos franceses e $x$, $y$, $z$
para variar sobre cidadãos alemães, então
\[
\lforall[a][\lforall x][(\Atom{\Obj{MarriedTo}}{a,x} \lif
(\Atom{\Obj{DrinksWine}}{a} \lor \lnot \Atom{\Obj{EatsWurst}}{x})]]
\]
afirma que, se algum francês for casado com um alemão, ou o
francês bebe vinho ou o alemão não come wurst.

A lógica de múltiplas ordenações pode ser embutida na lógica de primeira ordem de
maneira natural, reunindo todos os objetos dos !!{domain}s de múltiplas ordenações
em um único !!{domain} de primeira ordem, usando !!{predicate}s unários para acompanhar
as ordenações e relativizando quantificadores. Por exemplo, a linguagem de primeira ordem
correspondente ao exemplo acima teria !!{predicate}s unários
``$\Obj{German}$'' e ``$\Obj{French}$'', além das outras relações
descritas, com os requisitos de ordenação eliminados. Um quantificador ordenado $\lforall[x][!A]$, onde $x$ é !!a{variable} da ordenação alemã, traduz-se para
\[
\lforall[x][(\Atom{\Obj{German}}{x} \lif !A)].
\]
Precisamos adicionar axiomas que garantam que as ordenações sejam separadas---por exemplo,
$\lforall[x][\lnot (\Atom{\Obj{German}}{x} \land
  \Atom{\Obj{French}}{x})]$---, bem como axiomas que garantam que
``bebe vinho'' só vale de objetos que satisfazem o !!{predicate}
$\Atom{\Obj{French}}{x}$, etc. Com essas convenções e axiomas, não
é difícil mostrar que !!{sentence}s de múltiplas ordenações traduzem-se em
!!{sentence}s de primeira ordem, e !!{derivation}s de múltiplas ordenações traduzem-se
em !!{derivation}s de primeira ordem. Além disso, !!{structure}s de múltiplas ordenações
``traduzem-se'' em !!{structure}s de primeira ordem correspondentes e
vice-versa, de modo que também temos um teorema de completude para a lógica de múltiplas ordenações.

\end{document}
